\documentclass[10pt,a4paper]{amsart}
\usepackage{amsmath,amssymb,amsthm,graphicx}
\usepackage{verbatim}
\newenvironment{reference}{\comment}{\endcomment}
%\newenvironment{reference}{}{}
\newenvironment{slogan}{\comment}{\endcomment}
\newenvironment{history}{\comment}{\endcomment}

% For commutative diagrams we use Xy-pic
\usepackage[all]{xy}

% We use 2cell for 2-commutative diagrams.
\xyoption{2cell}
% Theorem environments.
%
\theoremstyle{plain}
\newtheorem{theorem}[subsection]{Theorem}
\newtheorem{proposition}[subsection]{Proposition}
\newtheorem{lemma}[subsection]{Lemma}

\theoremstyle{definition}
\newtheorem{definition}[subsection]{Definition}
\newtheorem{example}[subsection]{Example}
\newtheorem{exercise}[subsection]{Exercise}
\newtheorem{situation}[subsection]{Situation}

\theoremstyle{remark}
\newtheorem{remark}[subsection]{Remark}
\newtheorem{remarks}[subsection]{Remarks}

\numberwithin{equation}{subsection}

% Macros
%
\def\lim{\mathop{\mathrm{lim}}\nolimits}
\def\colim{\mathop{\mathrm{colim}}\nolimits}
\def\Spec{\mathop{\mathrm{Spec}}}
\def\Hom{\mathop{\mathrm{Hom}}\nolimits}
\def\Ext{\mathop{\mathrm{Ext}}\nolimits}
\def\SheafHom{\mathop{\mathcal{H}\!\mathit{om}}\nolimits}
\def\SheafExt{\mathop{\mathcal{E}\!\mathit{xt}}\nolimits}
\def\Sch{\mathit{Sch}}
\def\Mor{\mathop{\mathrm{Mor}}\nolimits}
\def\Ob{\mathop{\mathrm{Ob}}\nolimits}
\def\Sh{\mathop{\mathit{Sh}}\nolimits}
\def\NL{\mathop{N\!L}\nolimits}
\def\CH{\mathop{\mathrm{CH}}\nolimits}
\def\proetale{{pro\text{-}\acute{e}tale}}
\def\etale{{\acute{e}tale}}
\def\QCoh{\mathit{QCoh}}
\def\Ker{\mathop{\mathrm{Ker}}}
\def\Im{\mathop{\mathrm{Im}}}
\def\Coker{\mathop{\mathrm{Coker}}}
\def\Coim{\mathop{\mathrm{Coim}}}

% Boxtimes
%
\DeclareMathSymbol{\boxtimes}{\mathbin}{AMSa}{"02}

%
% Macros for moduli stacks/spaces
%
\def\QCohstack{\mathcal{QC}\!\mathit{oh}}
\def\Cohstack{\mathcal{C}\!\mathit{oh}}
\def\Spacesstack{\mathcal{S}\!\mathit{paces}}
\def\Quotfunctor{\mathrm{Quot}}
\def\Hilbfunctor{\mathrm{Hilb}}
\def\Curvesstack{\mathcal{C}\!\mathit{urves}}
\def\Polarizedstack{\mathcal{P}\!\mathit{olarized}}
\def\Complexesstack{\mathcal{C}\!\mathit{omplexes}}
% \Pic is the operator that assigns to X its picard group, usage \Pic(X)
% \Picardstack_{X/B} denotes the Picard stack of X over B
% \Picardfunctor_{X/B} denotes the Picard functor of X over B
\def\Pic{\mathop{\mathrm{Pic}}\nolimits}
\def\Picardstack{\mathcal{P}\!\mathit{ic}}
\def\Picardfunctor{\mathrm{Pic}}
\def\Deformationcategory{\mathcal{D}\!\mathit{ef}}
\usepackage{cleveref}
\crefname{originalsection}{original section}{original sections}
\Crefname{originalsection}{Original section}{Original sections}
\makeatletter
\newlabel{section-Riemann-Roch}{{0B5B}{source}{}{sourcepointer-1}{}}
\newlabel{section-Riemann-Roch@cref}{{[originalsection][0][]0B5B}{[1][0][]source}{}{}{}}
\makeatother
\crefname{originalreference}{original reference}{original references}
\Crefname{originalreference}{Original reference}{Original references}
\makeatletter
\newlabel{situation-Cohen-Macaulay-curve}{{0B5D}{source}{}{sourcepointer-2}{}}
\newlabel{situation-Cohen-Macaulay-curve@cref}{{[originalreference][0][]0B5D}{[1][0][]source}{}{}{}}
\makeatother
\title{Vanishing of H\textasciicircum{}1 twisted dualizing module for globally generated line bundles}
\author{The Stacks Project Authors}
\date{stacks-project:curves.tex:local-snapshot-20261012; GNU FDL 1.2 or later}
\begin{document}
\maketitle
\noindent\textit{Editorial note.} All delivered mathematical text is the author's own, copied from the pinned Stacks Project chapter file sources/stacks-curves/pinned-original-source.tex; no statement or proof step was rewritten, reordered or removed.\par\medskip
\noindent\textit{Editorial note.} The counted proof is the author's complete argument: two steps are implicit in the source and their mechanisms are recorded here: (a) the existence of s with st nonzero is the generator argument - t is a map L -> O\_X and vanishing on a generating family of global sections would force t = 0; (b) a nowhere-vanishing global section of an invertible sheaf trivializes it, so L cong O\_X; the duality input is the in-source Riemann-Roch section (navigation stub).\par\medskip
\noindent\textit{Editorial note.} The label lemma-vanishing-twist names the counted statement, and the proof environment immediately following it in the source is delivered verbatim as its proof.\par\medskip
\noindent\textit{Editorial note.} The counted statement is carried by the author's own lemma environment in the chapter source.\par\medskip
\noindent\textit{Allowed standard prerequisites (editorial).}\begin{itemize}
\item Serre duality on the Cohen-Macaulay curve. Same unit, Situation situation-Cohen-Macaulay-curve and the duality of Section section-Riemann-Roch. The proof reduces H\textasciicircum{}1(X, omega\_X tensor L) = 0 to H\textasciicircum{}0(X, L\textasciicircum{}\{tensor -1\}) = 0 by the duality discussed in the Riemann-Roch section.
\end{itemize}\medskip
\ifdefined\hypertarget\hypertarget{sourcepointer-1}{}\fi
\noindent\textit{Source reference.} Original section 0B5B: sources/stacks-curves/pinned-original-source.tex, source lines 1109--1109. Prerequisite pointer invoked at proof line(s) 1362; it names a standard fact of the same chapter that the counted argument uses, and every step of the counted proof built on it is retained verbatim in the delivered body.\par\medskip
\ifdefined\hypertarget\hypertarget{sourcepointer-2}{}\fi
\noindent\textit{Source reference.} Original reference 0B5D: sources/stacks-curves/pinned-original-source.tex, source lines 1307--1307. Prerequisite pointer invoked at proof line(s) 1355; it names a standard fact of the same chapter that the counted argument uses, and every step of the counted proof built on it is retained verbatim in the delivered body.\par\medskip
\begin{lemma}
\label{lemma-vanishing-twist}
In Situation \ref{situation-Cohen-Macaulay-curve}. Let
$\mathcal{L}$ be an invertible $\mathcal{O}_X$-module which is
globally generated and not isomorphic to $\mathcal{O}_X$. Then
$H^1(X, \omega_X \otimes \mathcal{L}) = 0$.
\end{lemma}

\begin{proof}
By duality as discussed in Section \ref{section-Riemann-Roch} we have to
show that $H^0(X, \mathcal{L}^{\otimes - 1}) = 0$. If not, then we can
choose a global section $t$ of $\mathcal{L}^{\otimes - 1}$
and a global section $s$ of $\mathcal{L}$ such that $st \not = 0$.
However, then $st$ is a constant multiple of $1$, by our assumption
that $H^0(X, \mathcal{O}_X) = k$. It follows that
$\mathcal{L} \cong \mathcal{O}_X$, which is a contradiction.
\end{proof}

\end{document}
